\documentclass[10pt,a4paper]{amsart}
\usepackage[margin=19mm]{geometry}
\usepackage{amsmath,amssymb,mathtools}
\usepackage[scr=rsfs,cal=euler]{mathalfa}
\usepackage{xfrac}
\usepackage[unicode,hidelinks,bookmarks=false]{hyperref}
\newcommand{\ie}{i.e.\ }
\newcommand{\cf}{cf.\ }
\newcommand{\resp}{respectively\ }
\newcommand{\Subsection}{\subsection}
\providecommand{\nobreakdash}{\nobreak\hbox{-}}
\setlength{\emergencystretch}{2em}
\multlinegap0pt
\usepackage{tikz}
\usetikzlibrary{cd}
\usepackage{amssymb}
\newtheorem{proposition}{Proposition}
\def\R{{\mathbb R}}%
\def\fv{\mathfrak{v}}
\title{The sub-Riemannian X-ray transform on $H$-type groups: Fourier slice theorems, injectivity sets and frequency support}
\author{Steven Flynn}
\date{Source: arXiv:2312.00594v2 (CC BY 4.0)}
\begin{document}
\maketitle

\noindent\emph{Source and licence.} The sub-Riemannian X-ray transform on $H$-type groups: Fourier slice theorems, injectivity sets and frequency support. Version arXiv:2312.00594v2 (CC BY 4.0), distributed by arXiv under the Creative Commons Attribution 4.0 International licence, http://creativecommons.org/licenses/by/4.0/, as recorded in the arXiv licence metadata for this exact version and shown on the abstract page https://arxiv.org/abs/2312.00594v2. Full licence notice: \texttt{licenses/arxiv-2312.00594-CC-BY-4.0.txt}.\par\smallskip

\noindent\textit{Editorial scope.} Counted unit: the article's proposition factorization (source0.tex lines 650--665). The article's complete original proof of it is delivered with it (source0.tex lines 666--701, one \texttt{proof} environment of 36 lines). No mathematics is written, repaired or re-derived: the statement and the proof are the article's own lines, in the article's own notation, and no retained line is substituted or re-flowed.  Retained uncounted as the source-local context the proof needs: lines 246--248.  Nothing was omitted from any retained span, so the package carries no omission record.  No retained line contains a citation, so the wrapper carries no bibliography and no bibliography entry.  The retained lines contain the article's own diagram code, which the wrapper typesets with the standard tikz packages; no image file is delivered and the package declares no asset dependency.  Editorial formatting only, in addition: an AMS \texttt{amsart} preamble that restates the article's own declarations and macros used by the retained lines, namely the environments \texttt{proposition} on separate counters, together with the standard packages those lines need.  The retained environments print in this document as Proposition 1 in order of appearance, whereas the article prints the same items with the numbering of its own full document.  The complete native source of the article as distributed in the arXiv e-print for this exact version is bundled unchanged as \texttt{sources/150831/source0.tex}.
\begin{align}\label{e:helicalSym}
	\gamma_{\nu, \lambda}(s+2\pi R)=(0, \pi R^2\hat\lambda)\,\gamma_{\nu, \lambda}(s).
\end{align}

\begin{proposition}\label{factorization}
	For any $\lambda\in \mathfrak{z}^*\setminus \{0\}$ and $\nu\in S(\fv^*)$, the X-ray transform $I_{\nu, \lambda}: L^1(G)\to L^1(G_\lambda)$ is well-defined, bounded and factors in the following way:
	\[ 
	\begin{tikzcd}
		L^1(G) \arrow[d, "P_{\lambda}" left=1]  \arrow[r, "{I}_{\nu, \lambda}"]  &L^1(G_\lambda)\\
		L^1(G_{\lambda}) \arrow[ur, "\widetilde{I}_{\nu, \lambda}" below=10, pos=0.75]
	\end{tikzcd}
	\]
	where the maps which we call \textit{central periodization} and the \textit{holonomy transform} are given by 
	\begin{align*}
		P_{\lambda}f(x, u)=\sum_{k\in\mathbb{Z}}f(x, u+\pi k |\lambda|^{-2}\hat{\lambda}), 
		\quad \quad
		\widetilde{I}_{\nu, \lambda}g(x, u)=\int_0^{2\pi |\lambda|^{-1}}g\left( (x, u)\gamma_{\nu, \lambda}(s)\right)ds
	\end{align*}
	for $f\in L^1(G)$ and $g\in L^1(G_\lambda)$.
\end{proposition}

\begin{proof}
	By the dominated convergence theorem we have
	\begin{align*}
		\int_{G_{\lambda}}P_{\lambda}f(x, u)d(x, u)=\int_G f(x, u)d(x, u), \quad f\in C_c^\infty(G).
	\end{align*}
	In particular $$\|P_{\lambda}f\|_{L^1(G_{\lambda})}\leq \|f\|_{L^1(G)}.$$
	
	Now $\widetilde{I}_{\nu, \lambda}f$ is just a group convolution operator on $G_{\lambda}$ by a compactly supported convolution kernel. Hence it is bounded. In particular, 
	\begin{align*}
		|\widetilde{I}_{\nu, \lambda}f(x, u)|\leq \int_0^{2\pi |\lambda|^{-1}}|f\left((x, u)\gamma_{\nu, \lambda}(s)\right)|ds
	\end{align*}
	and 
	\begin{align*}
		\int_{G_\lambda}|\widetilde{I}_{\nu, \lambda} f(x, u)|d(x, u)
		&=\int_0^{2\pi |\lambda|^{-1}} \int_{G_\lambda}|f\left(x, u\right)|d(x, u)ds.
	\end{align*}
	Therefore
	\begin{align*}
		\|\widetilde{I}_{\nu, \lambda} f\|_{L^1(G_\lambda)}\leq 2\pi  |\lambda|^{-1}\|f\|_{L^1(G_{\lambda})}.
	\end{align*}
	The decomposition follows from the helical symmetry \eqref{e:helicalSym} of the geodesics. For $f\in C_c^\infty(G)$ 
	\begin{align*}
		\widetilde{I}_{\nu, \lambda}\circ P_{\lambda}f(x, u)
		&=\int_0^{2\pi  |\lambda|^{-1}}\sum_{k\in\mathbb{Z}}f\left((x, u+\pi k|\lambda|^{-2}\hat{\lambda})\gamma_{\nu, \lambda}(s)\right)ds\\
		&=\int_0^{2\pi  |\lambda|^{-1}}\sum_{k\in\mathbb{Z}}f\left((x, u)\gamma_{\nu, \lambda}(s+2\pi k|\lambda|^{-1})\right)ds	\\
		&=\sum_{k\in\mathbb{Z}}\int_{2\pi k |\lambda|^{-1}}^{2\pi (k+1) |\lambda|^{-1}}f\left((x, u)\gamma_{\nu, \lambda}(s)\right)ds	\\
		&=\int_\R f\left((x, u)\gamma_{\nu, \lambda}(s)\right)ds=I_{\nu, \lambda}f(x, u)
	\end{align*}
	where the second to last equality follows from the dominated convergence theorem. 
	
	By the previous two statements, we have
	\begin{align*}
		\| I_{\nu, \lambda} f\|_{L^1(G_\lambda)}\leq 2\pi |\lambda|^{-1}\|f\|_{L^1(G)}.
	\end{align*}
	In particular $I_{\nu, \lambda}$ extends to a bounded map from $L^1(G)$ to $L^1(G_\lambda)$. 
\end{proof}

\end{document}
